\chapter{Number theory}

\section{Modular arithmetic}
	\kactlimport{ModularArithmetic.h}
	\kactlimport{ModInverse.h}
	\kactlimport{ModPow.h}
	\kactlimport{ModLog.h}
	\kactlimport{ModSum.h}
	\kactlimport{ModMulLL.h}
	\kactlimport{ModSqrt.h}

\section{Primality}
	\kactlimport{FastEratosthenes.h}
	\kactlimport{LinearSieve.h}
	\kactlimport{PrimeCounting.h}
	\kactlimport{MillerRabin.h}
	\kactlimport{Factor.h}

\section{Divisibility}
	\kactlimport{euclid.h}
	% \kactlimport{Euclid.java}
	\kactlimport{CRT.h}

	\subsection{Đẳng thức Bézout}
	Với $a \neq 0$, $b \neq 0$, đặt $d=\gcd(a,b)$.
	Khi đó, $d$ là số nguyên dương nhỏ nhất sao cho tồn tại các số nguyên
	$x,y$ thỏa mãn
	$$ax+by=d$$
	Nếu $(x,y)$ là một nghiệm thì mọi nghiệm đều có dạng
	$$\left(x+\frac{kb}{\gcd(a,b)}, y-\frac{ka}{\gcd(a,b)}\right),
	\quad k\in\mathbb{Z}$$

	% \kactlimport{phiFunction.h}

\section{Combinatorics modulo M}
    \kactlimport{Combinatorics.h}
    \kactlimport{nCrAny.h}

\section{Floor blocks}
    \kactlimport{FloorBlocks.h}

\section{Fractions}
	\kactlimport{ContinuedFractions.h}
	\kactlimport{FracBinarySearch.h}

\section{Bộ ba Pythagoras}
 Mọi bộ ba Pythagoras đều được sinh duy nhất bởi
 \[ a=k\cdot (m^{2}-n^{2}),\ \,b=k\cdot (2mn),\ \,c=k\cdot (m^{2}+n^{2}), \]
 với $m > n > 0$, $k > 0$, $m \bot n$ (nguyên tố cùng nhau), và một trong hai số $m, n$ chẵn.

\section{Số nguyên tố}
	$p=962592769$ thỏa $2^{21} \mid p-1$, có thể hữu ích (ví dụ cho NTT). Để băm,
	dùng 970592641 (31 bit), 31443539979727 (45 bit), 3006703054056749
	(52 bit). Có 78498 số nguyên tố nhỏ hơn $1\,000\,000$.

	Căn nguyên thủy tồn tại modulo mọi lũy thừa nguyên tố $p^a$, trừ $p = 2, a > 2$, và có đúng $\phi(\phi(p^a))$ căn nguyên thủy.
	Với $p = 2, a > 2$, nhóm $\mathbb Z_{2^a}^\times$ đẳng cấu với $\mathbb Z_2 \times \mathbb Z_{2^{a-2}}$.

\section{Ước lượng}
	$\sum_{d|n} d = O(n \log \log n)$.

	Số ước của $n$ nhiều nhất khoảng 100 khi $n < 5\cdot 10^4$, 500 khi $n < 10^7$, 2000 khi $n < 10^{10}$, 200\,000 khi $n < 10^{19}$.

\section{Hàm Mobius}
\[
	\mu(n) = \begin{cases} 0 & n \textrm{ có ước chính phương}\\ 1 & n \textrm{ có số lượng ước nguyên tố chẵn}\\ -1 & n \textrm{ có số lượng ước nguyên tố lẻ}\\\end{cases}
\]
  (Số ước nguyên tố tính theo các ước nguyên tố phân biệt, $n$ không có ước chính phương.)

  Công thức nghịch đảo Mobius:
  \[ g(n) = \sum_{d|n} f(d) \Leftrightarrow f(n) = \sum_{d|n} \mu(d)g(n/d) \]
  Một số công thức hữu ích khác:

  $ \sum_{d | n} \mu(d) = [ n = 1] $ (rất hay dùng)

  $ g(n) = \sum_{n|d} f(d) \Leftrightarrow f(n) = \sum_{n|d} \mu(d/n)g(d)$

 $ g(n) = \sum_{1 \leq m \leq n} f(\left\lfloor\frac{n}{m}\right \rfloor ) \Leftrightarrow f(n) = \sum_{1\leq m\leq n} \mu(m)g(\left\lfloor\frac{n}{m}\right\rfloor)$
